\chapter{Pendahuluan}

\section{Latar Belakang}
\TemplateTodo{Jelaskan konteks umum topik, kondisi ideal, masalah nyata, kesenjangan penelitian atau kebutuhan praktis, serta alasan penelitian perlu dilakukan.}

Tuliskan latar belakang secara runtut dari gambaran besar menuju masalah khusus \DocTypeNameLower. Gunakan sumber ilmiah yang dapat diverifikasi, misalnya artikel jurnal, prosiding, buku, standar, atau dokumen resmi \cite{exampleReference}.

\TemplateTip{Hindari langsung membahas detail teknis tanpa menjelaskan masalah. Setiap klaim penting sebaiknya didukung sitasi.}

\section{Rumusan Masalah}
\TemplateTodo{Ubah masalah penelitian menjadi pertanyaan yang spesifik, terukur, dan dapat dijawab melalui metodologi.}

\begin{enumerate}
  \item Bagaimana merancang pendekatan untuk menyelesaikan masalah utama pada objek penelitian?
  \item Bagaimana mengevaluasi kinerja pendekatan yang diusulkan berdasarkan metrik yang sesuai?
  \item Faktor apa saja yang memengaruhi hasil penelitian berdasarkan analisis yang dilakukan?
\end{enumerate}

\section{Batasan Masalah}
\TemplateTodo{Tuliskan ruang lingkup agar penelitian tidak melebar. Batasan dapat berupa data, metode, lokasi, periode, variabel, alat, atau asumsi.}

\section{Tujuan Penelitian}
\TemplateTodo{Selaraskan tujuan dengan rumusan masalah. Gunakan kata kerja yang dapat diverifikasi seperti menganalisis, merancang, mengimplementasikan, mengevaluasi, atau membandingkan.}

\section{Manfaat Penelitian}
\TemplateTodo{Jelaskan manfaat teoretis dan praktis. Hindari klaim terlalu besar; tuliskan kontribusi yang realistis sesuai skala \DocTypeNameLower.}
